\documentclass[10pt,a4paper]{amsart}
\usepackage{amsmath,amssymb,amsthm,graphicx}
% Retained minimal preamble for the Stacks Project chapter flat.tex (More on Flatness).
% Verbatim excerpt of the chapter preamble: Xy-pic setup (native preamble lines 14-18) and the native
% theorem declarations with the mathematical macros (native preamble lines 154-217). Copied mechanically;
% nothing added, removed or redefined. The wrapper class and the amsmath/amssymb/amsthm/graphicx base
% packages are supplied by the assembler, so the class/externaldocument part of the native preamble is omitted.
% For commutative diagrams we use Xy-pic
\usepackage[all]{xy}

% We use 2cell for 2-commutative diagrams.
\xyoption{2cell}
% Theorem environments.
%
\theoremstyle{plain}
\newtheorem{theorem}[subsection]{Theorem}
\newtheorem{proposition}[subsection]{Proposition}
\newtheorem{lemma}[subsection]{Lemma}

\theoremstyle{definition}
\newtheorem{definition}[subsection]{Definition}
\newtheorem{example}[subsection]{Example}
\newtheorem{exercise}[subsection]{Exercise}
\newtheorem{situation}[subsection]{Situation}

\theoremstyle{remark}
\newtheorem{remark}[subsection]{Remark}
\newtheorem{remarks}[subsection]{Remarks}

\numberwithin{equation}{subsection}

% Macros
%
\def\lim{\mathop{\mathrm{lim}}\nolimits}
\def\colim{\mathop{\mathrm{colim}}\nolimits}
\def\Spec{\mathop{\mathrm{Spec}}}
\def\Hom{\mathop{\mathrm{Hom}}\nolimits}
\def\Ext{\mathop{\mathrm{Ext}}\nolimits}
\def\SheafHom{\mathop{\mathcal{H}\!\mathit{om}}\nolimits}
\def\SheafExt{\mathop{\mathcal{E}\!\mathit{xt}}\nolimits}
\def\Sch{\mathit{Sch}}
\def\Mor{\mathop{\mathrm{Mor}}\nolimits}
\def\Ob{\mathop{\mathrm{Ob}}\nolimits}
\def\Sh{\mathop{\mathit{Sh}}\nolimits}
\def\NL{\mathop{N\!L}\nolimits}
\def\CH{\mathop{\mathrm{CH}}\nolimits}
\def\proetale{{pro\text{-}\acute{e}tale}}
\def\etale{{\acute{e}tale}}
\def\QCoh{\mathit{QCoh}}
\def\Ker{\mathop{\mathrm{Ker}}}
\def\Im{\mathop{\mathrm{Im}}}
\def\Coker{\mathop{\mathrm{Coker}}}
\def\Coim{\mathop{\mathrm{Coim}}}

% Boxtimes
%
\DeclareMathSymbol{\boxtimes}{\mathbin}{AMSa}{"02}

%
% Macros for moduli stacks/spaces
%
\def\QCohstack{\mathcal{QC}\!\mathit{oh}}
\def\Cohstack{\mathcal{C}\!\mathit{oh}}
\def\Spacesstack{\mathcal{S}\!\mathit{paces}}
\def\Quotfunctor{\mathrm{Quot}}
\def\Hilbfunctor{\mathrm{Hilb}}
\def\Curvesstack{\mathcal{C}\!\mathit{urves}}
\def\Polarizedstack{\mathcal{P}\!\mathit{olarized}}
\def\Complexesstack{\mathcal{C}\!\mathit{omplexes}}
% \Pic is the operator that assigns to X its picard group, usage \Pic(X)
% \Picardstack_{X/B} denotes the Picard stack of X over B
% \Picardfunctor_{X/B} denotes the Picard functor of X over B
\def\Pic{\mathop{\mathrm{Pic}}\nolimits}
\def\Picardstack{\mathcal{P}\!\mathit{ic}}
\def\Picardfunctor{\mathrm{Pic}}
\def\Deformationcategory{\mathcal{D}\!\mathit{ef}}

% --- appended 2026-10-12 (Chief): the project own definitions of the reference/slogan/history
% display environments, copied mechanically from the shared Stacks preamble lines 8-12
% (\usepackage{verbatim} plus the three \newenvironment lines, which the chapter preamble keeps
% verbatim from the project preamble). Nothing added, removed or redefined: these environments
% carry no mathematical content and the project itself typesets them as comments. Without these
% three lines a counted statement whose own text uses them fails with Environment slogan undefined.
\usepackage{verbatim}
\newenvironment{reference}{\comment}{\endcomment}
\newenvironment{slogan}{\comment}{\endcomment}
\newenvironment{history}{\comment}{\endcomment}
\usepackage{cleveref}
\crefname{originaldefinition}{original definition}{original definitions}
\Crefname{originaldefinition}{Original definition}{Original definitions}
\makeatletter
\newlabel{definition-big-small-h}{{0EU1}{source}{}{sourcepointer-1}{}}
\newlabel{definition-big-small-h@cref}{{[originaldefinition][0][]0EU1}{[1][0][]source}{}{}{}}
\makeatother
\crefname{originallemma}{original lemma}{original lemmas}
\Crefname{originallemma}{Original lemma}{Original lemmas}
\makeatletter
\newlabel{lemma-refine-check-h}{{0EVI}{source}{}{sourcepointer-2}{}}
\newlabel{lemma-refine-check-h@cref}{{[originallemma][0][]0EVI}{[1][0][]source}{}{}{}}
\makeatother
\crefname{originallemma}{original lemma}{original lemmas}
\Crefname{originallemma}{Original lemma}{Original lemmas}
\makeatletter
\newlabel{descent-lemma-sheaf-condition-holds}{{03DT}{source}{}{sourcepointer-3}{}}
\newlabel{descent-lemma-sheaf-condition-holds@cref}{{[originallemma][0][]03DT}{[1][0][]source}{}{}{}}
\makeatother
\crefname{originalsection}{original section}{original sections}
\Crefname{originalsection}{Original section}{Original sections}
\makeatletter
\newlabel{example-one-generator}{{0EVG}{source}{}{sourcepointer-4}{}}
\newlabel{example-one-generator@cref}{{[originalsection][0][]0EVG}{[1][0][]source}{}{}{}}
\makeatother
\crefname{originalsection}{original section}{original sections}
\Crefname{originalsection}{Original section}{Original sections}
\makeatletter
\newlabel{example-two-generators}{{0EVH}{source}{}{sourcepointer-5}{}}
\newlabel{example-two-generators@cref}{{[originalsection][0][]0EVH}{[1][0][]source}{}{}{}}
\makeatother
\crefname{originallemma}{original lemma}{original lemmas}
\Crefname{originallemma}{Original lemma}{Original lemmas}
\makeatletter
\newlabel{lemma-funny-blow-up}{{0EVJ}{source}{}{sourcepointer-6}{}}
\newlabel{lemma-funny-blow-up@cref}{{[originallemma][0][]0EVJ}{[1][0][]source}{}{}{}}
\makeatother
\title{The limit of the structure sheaf along the Frobenius is a sheaf for the h topology}
\author{The Stacks Project Authors}
\date{stacks-project:flat.tex:local-snapshot-20261011; GNU FDL 1.2 or later}
\begin{document}
\maketitle
\noindent\textit{Editorial note.} All delivered mathematical text is the author's own, copied from the pinned Stacks Project chapter file flat.tex; no statement or proof step was rewritten, reordered or removed.\par\medskip
\noindent\textit{Editorial note.} The preamble is the chapter's own retained minimal preamble (sources/stacks-flat/retained-preamble.tex), a verbatim excerpt of the native Xy-pic setup and of the native theorem declarations and mathematical macros, followed by the project own definitions of the reference, slogan and history display environments (verbatim package plus three newenvironment lines) that the retained file carries at its end and that the project itself typesets as comments. No macro was invented or redefined, and the wrapper class and the amsmath/amssymb/amsthm/graphicx base packages are supplied by the assembler.\par\medskip
\noindent\textit{Editorial note.} The counted statement and proof contain no cross reference, no citation and no included file, so the delivered body is self contained.\par\medskip
\noindent\textit{Editorial note.} Canonical location of the counted proposition: Stacks tag 0EVL (sources/stacks-flat/pinned-original-source.tex line 11931).\par\medskip
\noindent\textit{Allowed standard prerequisites (editorial).}\begin{itemize}
\item The four conditions characterising h sheaves. Same chapter, Lemma labelled lemma-refine-check-h, cited at proof lines 11941-11943. Provides the checklist verified in the proof.
\item The fppf sheaf property of the structure sheaf. Descent, Lemma labelled descent-lemma-sheaf-condition-holds, cited at proof lines 11945-11947. Settles condition (2) together with commutation of the limit with equalizers.
\item The one generator example and the nilpotence of fA cap J. Same chapter, Example labelled example-one-generator, cited at proof lines 11953-11955. Supplies the algebra setting of condition (3).
\item The two generators example and the funny blowup. Same chapter, Example labelled example-two-generators, invoked at proof lines 11969-11971, and Lemma labelled lemma-funny-blow-up, cited at proof lines 11969-11973. Settles condition (4).
\end{itemize}\medskip
\ifdefined\hypertarget\hypertarget{sourcepointer-1}{}\fi
\noindent\textit{Source reference.} Original definition 0EU1: sources/stacks-flat/pinned-original-source.tex, source lines 10504--10504. Same-chapter pointer into the author environment definition at source line 10504; invoked at proof line None of the counted proof. It is retained only as a navigation pointer; every load-bearing step is inside the delivered span.\par\medskip
\ifdefined\hypertarget\hypertarget{sourcepointer-2}{}\fi
\noindent\textit{Source reference.} Original lemma 0EVI: sources/stacks-flat/pinned-original-source.tex, source lines 11501--11501. Same-chapter pointer into the author environment lemma at source line 11501; invoked at proof line 11943 of the counted proof. It is retained only as a navigation pointer; every load-bearing step is inside the delivered span.\par\medskip
\ifdefined\hypertarget\hypertarget{sourcepointer-3}{}\fi
\noindent\textit{Source reference.} Original lemma 03DT: sources/stacks-chow/tags.txt, source lines 4432--4432. Cross-chapter standard prerequisite located in the project tag list staged in the local snapshot; invoked at proof line 11947 of the counted proof. It is retained only as a navigation pointer; every load-bearing step is inside the delivered span.\par\medskip
\ifdefined\hypertarget\hypertarget{sourcepointer-4}{}\fi
\noindent\textit{Source reference.} Original section 0EVG: sources/stacks-flat/pinned-original-source.tex, source lines 11474--11474. Same-chapter pointer into the author environment example at source line 11474; invoked at proof line 11953 of the counted proof. It is retained only as a navigation pointer; every load-bearing step is inside the delivered span.\par\medskip
\ifdefined\hypertarget\hypertarget{sourcepointer-5}{}\fi
\noindent\textit{Source reference.} Original section 0EVH: sources/stacks-flat/pinned-original-source.tex, source lines 11488--11488. Same-chapter pointer into the author environment example at source line 11488; invoked at proof line 11972 of the counted proof. It is retained only as a navigation pointer; every load-bearing step is inside the delivered span.\par\medskip
\ifdefined\hypertarget\hypertarget{sourcepointer-6}{}\fi
\noindent\textit{Source reference.} Original lemma 0EVJ: sources/stacks-flat/pinned-original-source.tex, source lines 11782--11782. Same-chapter pointer into the author environment lemma at source line 11782; invoked at proof line 11973 of the counted proof. It is retained only as a navigation pointer; every load-bearing step is inside the delivered span.\par\medskip
\begin{lemma}
\label{lemma-h-sheaf-lim-F}
Let $p$ be a prime number. Let $S$ be a scheme over $\mathbf{F}_p$.
Let $(\Sch/S)_h$ be a site as in Definition \ref{definition-big-small-h}.
The rule
$$
\mathcal{F}(X) = \lim_F \Gamma(X, \mathcal{O}_X)
$$
defines a sheaf on $(\Sch/S)_h$.
\end{lemma}

\begin{proof}
To prove $\mathcal{F}$ is a sheaf, let's check conditions
(1), (2), (3), and (4) of Lemma \ref{lemma-refine-check-h}.
Condition (1) holds because limits of sheaves are sheaves
and $\mathcal{O}$ is a Zariski sheaf. Condition (2) holds because
$\mathcal{O}$ is an fppf sheaf (Descent, Lemma
\ref{descent-lemma-sheaf-condition-holds}) and if $A$ is the equalizer
of two maps $B \to C$ of $\mathbf{F}_p$-algebras, then $\lim_F A$
is the equalizer of the two maps $\lim_F B \to \lim_F C$.

\medskip\noindent
We check condition (3). Let $A, f, J$ be as in
Example \ref{example-one-generator}.
We have to show that
\begin{align*}
\lim_F A
& \to
\lim_F A/J \times_{\lim_F A/fA + J} \lim_F A/fA \\
& =
\lim_F (A/J \times_{A/fA + J} A/fA) \\
& =
\lim_F A/(fA \cap J)
\end{align*}
is bijective. Since $J$ is annihilated by a power of $f$ we see that
$\mathfrak a = fA \cap J$ is a nilpotent ideal, i.e., there exists an
$n$ such that $\mathfrak a^n = 0$. It is straightforward
to verify that in this case $\lim_F A \to \lim_F A/\mathfrak a$
is bijective.

\medskip\noindent
We check condition (4). Let $X, X', Z, E$ be as in
Example \ref{example-two-generators}. By
Lemma \ref{lemma-funny-blow-up} and the same argument as above
we see that
$$
\mathcal{F}(X) = \lim_F \Gamma(X, \mathcal{O}_X)
\longrightarrow
\lim_F \Gamma(X', \mathcal{O}_{X'}) = \mathcal{F}(X')
$$
is bijective. Since $E = \mathbf{P}^1_Z$ in this case we also
see that $\mathcal{F}(Z) \to \mathcal{F}(E)$ is bijective.
Thus the conclusion holds in this case as well.
\end{proof}

\end{document}
